\section{From CLEWS to OG-Core}
\label{sec:linkages}

\subsection{Electricity cost}

The preferred measure of electricity cost would be the cost of supplying one additional
unit. The current CLEWS runs do not produce a usable series because the relevant supply
constraint binds only in scattered years. The link therefore uses a curated
electricity-cost index where one is available and otherwise reconstructs levelized cost
from generation, capital, fuel, and operating costs. It passes the reform-to-baseline
ratio, so the calculation does not depend on converting monetary levels between the
models.

This is an average-cost proxy, not an observed tariff or marginal price. That distinction
limits its interpretation but still provides a scenario-derived measure of how the cost
of the electricity system changes.

The link applies the ratio to two separate uses of electricity. For households, it changes
the price of final electricity consumption. Lower-income households can be more exposed
because electricity includes a minimum-consumption component in OG-Core. The price change
is represented through a consumption-price adjustment, but the resulting notional revenue
is returned to households as a lump-sum transfer: a higher system cost is not government
tax revenue.

For firms, the link uses the country's input--output accounts to identify each industry's direct
and indirect exposure to electricity costs. Electricity purchased by the electricity
industry itself is removed before the cost is propagated, preventing overlap between
self-use and sales to other industries. The result is an industry-specific cost increase
rather than a uniform shock to the economy.

The current firm-side representation is an approximation. OG-Core does not yet include
electricity as an explicit intermediate input, so firms cannot substitute away from it as
its price changes. Adding an explicit energy input would allow electricity demand, output,
and prices to adjust together inside the economic model.

\subsection{Public grid investment}

CLEWS reports additional investment in grids and transmission separately from investment
in generation. The link treats this as public investment. It raises the flow of public
capital in OG-Core, where infrastructure contributes to production and its financing
affects the government budget.

The financing assumption matters. Debt-financed investment competes for saving and can
displace private capital; tax financing and external finance create different effects.
The current mapping establishes this mechanism, but the grid-investment difference in the
test application is too small to validate a substantial economic response. The evidence
therefore supports the transformation and solve, not a calibrated investment magnitude.

\subsection{Generation capital intensity}

CLEWS also shows whether the reform changes the capital-cost share of the generation
fleet. The link uses that ratio to change the production structure of the electricity
industry in OG-Core: capital receives a different share relative to labour and other
costs.

This does not mechanically increase the amount of capital employed. A higher capital
share changes production cost, output, price, and revenue; the resulting capital stock
depends on all of them. Controlled runs confirm that the capital response can differ from
the change in the capital share. This relationship should therefore be read as a change
in how electricity is produced, not as an investment programme.

Private investment support is represented separately through the cost of capital. The
two adjustments can coexist only when they correspond to different source information:
one describing the technology mix, the other describing a policy incentive.

\subsection{Air pollution, survival, and labour}

The health relationship uses the change in CLEWS particulate emissions to adjust two
parts of OG-Core. The mortality effect changes age-specific survival and causes the model
to recompute the population. The morbidity effect changes effective labour among people
of working age. Cleaner air can therefore raise near-term productivity while also
changing the future balance between workers and retirees.

The age profiles come from burden-of-disease data. A calibrated dose-response translates
the power-sector emissions change into the share of the national health burden affected.
This avoids treating a percentage reduction in power-sector emissions as the same
percentage reduction in all pollution-related illness and death.

The distinction between sourced evidence and model calibration is important. The burden
and age profiles are external data. The emissions-to-burden translation and the effect of
morbidity on productivity remain calibrated assumptions. The relationship is implemented
and produces the intended demographic and labour responses, but its magnitude remains
uncertain.
